\documentclass{article}
\usepackage[utf8]{inputenc}
\usepackage{mathtools,amssymb}
\usepackage{lipsum}
\usepackage{hyperref}
\usepackage{layout}
\usepackage{geometry}
\usepackage{listings}
\usepackage{color}
\usepackage{xcolor}
\usepackage{amsthm}
\newcommand*{\vertbar}[1][10ex]{\rule[-1ex]{0.5pt}{#1}}
\newcommand*{\horzbar}{\rule[.5ex]{2.5ex}{0.5pt}}
\newcommand{\norm}[2]{\left\lVert#1\right\rVert_#2}
\newcommand{\abs}[1]{\lvert#1\rvert}

\title{6.2202 class notes}
\author{Erich Trieschman}
\date{2026 Fall semester}

% compact spacing
\usepackage{setspace}\setstretch{0.95}
\usepackage{parskip}\setlength{\parskip}{4pt}
\setlength{\parindent}{0pt}
\usepackage{enumitem}
\setlist[itemize]{topsep=2pt, itemsep=2pt, parsep=0pt, partopsep=0pt}
\setlist[enumerate]{topsep=2pt, itemsep=2pt, parsep=0pt, partopsep=0pt}
\usepackage{titlesec}
\titlespacing*{\section}{0pt}{1.5ex plus 0.3ex minus 0.3ex}{0.5ex}
\titlespacing*{\subsection}{0pt}{2.5ex plus 0.2ex minus 0.2ex}{0.3ex}
\titlespacing*{\subsubsection}{0pt}{1.5ex plus 0.2ex minus 0.2ex}{0.2ex}
\titlespacing*{\paragraph}{0pt}{1.0ex plus 0.1ex minus 0.1ex}{1.0ex}
\definecolor{secblue}{RGB}{31,78,140}
\definecolor{parorange}{RGB}{204,102,0}
\titleformat{\section}{\normalfont\bfseries\color{secblue}}{\thesection}{1.0em}{}
\titleformat{\subsection}{\small\bfseries\color{secblue}}{\thesubsection}{0.5em}{}
\titleformat{\paragraph}[runin]{\small\bfseries\color{parorange}}{\theparagraph}{}{}{}

\newcommand{\userMarginInMm}{10}
\geometry{
%  legal,
 left=\userMarginInMm mm,
 right=\userMarginInMm mm,
 top=\userMarginInMm mm,
 bottom=\userMarginInMm mm,
 footskip=5mm}


 \hypersetup{
    colorlinks=true,
    linkcolor=blue,
    filecolor=magenta,      
    urlcolor=cyan,
    pdfpagemode=FullScreen,
}

\definecolor{dkgreen}{rgb}{0,0.6,0}
\definecolor{gray}{rgb}{0.5,0.5,0.5}
\definecolor{mauve}{rgb}{0.58,0,0.82}
\newcommand{\todo}[1]{\textit{\textcolor{orange}{ [#1]}}}
    
\lstset{frame=tb, language=Python,
    aboveskip=3mm,
  belowskip=3mm,
  showstringspaces=false,
  columns=flexible,
  basicstyle={\small\ttfamily},
  numbers=none,
  numberstyle=\tiny\color{gray},
  keywordstyle=\color{blue},
  commentstyle=\color{dkgreen},
  stringstyle=\color{mauve},
  breaklines=true,
  breakatwhitespace=true,
  tabsize=3}

\begin{document}
% \maketitle
% \tableofcontents

\section{Circuit review}
% \paragraph{Laws.}
% KCL: at every node, $\sum(\text{currents in}) = \sum(\text{currents out})$.
% KVL: around every closed loop, $\sum V = 0$.

\paragraph{Phasors.}
$v(t) = V_{max}\cos(\omega t + \phi)$.
Positive $\phi$ shifts the curve left and waveform occurs earlier.
RMS value: $\tilde V = V_{rms} = V_{max}/\sqrt2$.
% Power engineering reports $V_{rms}$, not $V_{max}$; state which you use.
$v(t) = \text{Re}\{\tilde Ve^{\omega t + \phi}\}$.
Phasor drops $e^{\omega t}$: $V = \tilde V\angle\phi$.
$j = \sqrt{-1} = 1\angle90^\circ$, and $1/j = -j$.
% Rectangular $\iff$ polar: $a + jb = \sqrt{a^2 + b^2}\,\angle\operatorname{atan2}(b,a)$, and $r\angle\theta = r\cos\theta + jr\sin\theta = re^{j\theta}$.
% $(r_1\angle\theta_1)/(r_2\angle\theta_2) = (r_1/r_2)\angle(\theta_1 - \theta_2)$.
Conjugate: $(r\angle\theta)^* = r\angle(-\theta)$.

\paragraph{Elements.}
% Sign convention: $I$ is defined entering the $+$ terminal of passive $V$.
Impedance $Z = R + jX$ [$\Omega$], with resistance $R$ and reactance $X$.
Polar: $Z = \abs Z\angle\theta$, $\abs Z = \sqrt{R^2 + X^2}$, $\tan\theta = X/R$.
Admittance $Y = 1/Z = G + jB$ [S]
% , with conductance $G$ and susceptance $B$; in polar $Y = (1/\abs Z)\angle(-\theta)$.
Ohm's law: $V = IZ$, or $I = VY$.
With $V = \tilde V\angle\phi_V$ and $I = \tilde I\angle\phi_I$, $\theta = \phi_V - \phi_I$.
Inductor $L$: $v = L\,di/dt$, so $X_L = \omega L$.
Capacitor $C$: $i = C\,dv/dt$, so $X_C = -1/(\omega C)$.
Inductive (``lag''): $\theta > 0$, $\phi_I < \phi_V$, current occurs after voltage.
Capacitive (``lead''): $\theta < 0$, $\phi_I > \phi_V$, current occurs before voltage.

\paragraph{Solving.}
Series: $Z_{eq} = Z_1 + Z_2 + \dots$; same $I$.
Parallel: $Y_{eq} = Y_1 + Y_2 + \dots$, or $Z_{eq} = Z_1 Z_2/(Z_1 + Z_2)$ for two; same $V$.
Voltage divider (series): $V_2 = V\cdot Z_2/(Z_1 + Z_2)$.
Current divider (parallel): $I_2 = I\cdot Z_1/(Z_1 + Z_2)$.
% Thevenin: a source $V_{Th}$ in series with $Z_{Th}$.
% $V_{Th}$ is the open-circuit voltage at the terminals.
% $Z_{Th}$ is the impedance seen at the terminals with independent sources zeroed (voltage source $\to$ short, current source $\to$ open), or $Z_{Th} = V_{oc}/I_{sc}$.
Norton: a current source $I_N = V_{Th}/Z_{Th}$ in parallel with $Z_N = Z_{Th}$.
$v = V_{\rm TH} + R_{\rm TH}i$ with $i$ coming into the circuit.

\paragraph{Transients.}
RC or RL circuit with a DC source: $x(t\geq 0) = x(\infty) + [x(0^+) - x(\infty)]\,e^{-t/\tau}$.
Time constant $\tau = R_{Th}C$ (RC) or $L/R_{Th}$ (RL), with $R_{Th}$ seen by $v_C$ or $i_L$.
$v_C$ and $i_L$ cannot jump: $x(0^+) = x(0^-)$.
$x(\infty)$: capacitor open, inductor short.
% Any other variable at $0^+$: redraw with $C$ as a voltage source $v_C(0^+)$ (or $L$ as a current source $i_L(0^+)$) and solve.
Origin: particular solution (steady state) plus homogeneous $Ae^{-t/\tau}$, with $A$ set by $x(0^+)$.
% Power absorbed by any element: $p(t) = v(t)\,i(t)$.
Stored energy is $w(t) = \int_0^t p(t)\,dt = \int v(t)\,i(t)\,dt$.
Capacitor: $i(t)\,dt = C\,dv$, so $w(t) = \int v(t)\,C\,dv = \tfrac12 C v(t)^2$.
Inductor: $v(t)\,dt = L\,di$. %, so $w(t) = \int L\,i(t)\,di = \tfrac12 L i(t)^2$.



\section{Power}
\paragraph{Definitions.}
Complex power $S = VI^* = P[W] + jQ[VAr]$ [VA] (>0 absorbed).
% $S = P + jQ$, with real power $P$ [W] and reactive power $Q$ [VAr].
Polar: $S = \tilde V\tilde I\angle\theta$, where $\theta = \phi_V - \phi_I$.
So $P = \tfrac12 V_{max}I_{max}\cos\theta=\tilde V\tilde I\cos\theta$, $Q = P\tan\theta$.
Power factor $\mathrm{pf} = \cos\theta = P/\abs S$.
Lagging: $\theta > 0$, inductive, $Q > 0$, absorbs $Q$.
Leading: $\theta < 0$, capacitive, $Q < 0$, supplies $Q$.
Single element: $S = \tilde I^2 Z = \tilde V^2Y^*$.
$\theta = \pm\arccos(\mathrm{pf})$, positive if lagging.

% \paragraph{Load to element.}
% Given $\abs S$ (or $P$) and a pf, lagging or leading.
% $\theta = \pm\arccos(\mathrm{pf})$, positive if lagging.
% As a current: $I = (S/V)^*$, so $\tilde I = \abs S/\tilde V$ and $\phi_I = \phi_V - \theta$.
% As an impedance: $Z = \tilde V^2/S^*$.

% \paragraph{Conservation.}
% $\sum S = 0$ at a node and $P$ and $Q$ balance separately.
% Power-factor correction: given $P$, pf $\theta_1$, bus $V = \tilde V\angle0$, target $\theta_2$; find $C$ in parallel with the load.
% $S_L = P + jP\tan\theta_1$; $S_C = \abs V^2 Y^* = -j\omega C\tilde V^2$; $S_{src} = S_L + S_C = P + j(P\tan\theta_1 - \omega C\tilde V^2)$.
% Target: $Q_{src} = P\tan\theta_2$, since $Q = P\tan\theta$ and $P$ is unchanged.




\section{Three-phase and per unit}
\paragraph{Definitions.}
Balanced: phase b is phase a $\angle{-120^\circ}$, phase c is phase a $\angle{+120^\circ}$, for every voltage and current; solve phase a only.
% Balanced line currents sum to zero, so the neutral carries no current and may be omitted.
$V_{ln}$: phase to neutral, e.g.\ $V_{an}$.
$V_{ll}$: phase to phase, $V_{ab} = V_{an} - V_{bn} = \sqrt3\,V_{an}\angle{+30^\circ}$.
% Voltage ratings are $V_{ll}$ unless stated. 
Relationship holds for both connection types.
Line current: in a conductor between source and load ($I_a$).
Phase current: in one load element; equals the line current in Y ($I_{ab}$ in $\Delta$).

\paragraph{Connections.}
Y: element between phase and neutral; element voltage $V_{ln}$, element current $I_a = V_{an}/Z_Y$.
$\Delta$: element between two phases; element voltage $V_{ll}$, element current $I_{ab} = V_{ab}/Z_\Delta$.
$I_a = I_{ab} - I_{ca} = \sqrt3\,I_{ab}\angle{-30^\circ}$.
$\Delta$ to Y: $Z_Y = Z_\Delta/3$ gives the same line currents.

\paragraph{Solving.}
Convert every $\Delta$ to Y and every $V_{ll}$ to $V_{ln} = V_{ll}/\sqrt3\,\angle{-30^\circ}$.
Draw phase a: source $V_{an}$, series source and line impedances, Y-equivalent load, neutral return; shunt elements line to neutral.
A load given as $\abs{S_{3\phi}}$ (or $P$) and a pf at load voltage $V_{ll}$: $S_\phi = \abs{S_{3\phi}}\angle\theta/3 = P/(3\text{pf})\angle\theta$ with $\theta$ from pf; $Z_Y = \tilde V_{ln}^2/S_\phi^*$.

\paragraph{Power.}
$S_{3\phi} = 3V_{an}I_a^* = \sqrt3\,\tilde V_{ll}\tilde I_{line}\angle\theta$.
Y: $\abs S_{3\phi} = 3\tilde V_{ln}^2/\abs{Z_Y}$. $\Delta$: $\abs S_{3\phi} = 3\tilde V_{ll}^2/\abs{Z_\Delta} = 9\tilde V_{ln}^2/\abs{Z_\Delta}$.

\paragraph{Per unit.}
$S_b$ for the network, three-phase, typically 100 MVA; one $V_b$ per zone, line-to-line, nominal.
pu $=$ actual$/$base; scales magnitues only.
Single-phase $I_b = S_b/V_b$, $Z_b = V_b^2/S_b$.
Three-phase $V_{b,ln} = V_{b,ll}/\sqrt3$, $I_b = S_b/(\sqrt3 V_{b,ll})$, $Z_b = V_{b,ll}^2/S_b$; 
Base change: $Z_{pu,new} = Z_{pu,old}\,(V_{b,old}/V_{b,new})^2\,(S_{b,new}/S_{b,old})$.
In pu, $\abs{V_{ln}} = \abs{V_{ll}}$ and $S_\phi = S_{3\phi}$.
Across a transformer, set the other side's $V_b$ by its line-to-line ratio; $Z_{pu}$ is the same on both sides and the ideal transformer drops out.




\section{Power flow and admittance matrices}
\paragraph{Admittance matrix.}
Ground as reference; generators as current injections $I_k$; impedance loads and line charging as shunt admittances to ground; series impedances as $y_{km} = 1/Z_{km}$.
KCL at each node gives $I_k = \sum_m Y_{km}V_m\implies I = YV$: $Y_{kk}$ = sum of all admittances touching $k$; $Y_{km} = -y_{km}$,; symmetric, singular without a shunt.
Branch current $y_{km}(V_k - V_m)$.
$Y = A^\top\operatorname{diag}(y_{\rm series})A + \operatorname{diag}(y_{\rm shunt})$.

\paragraph{Power flow.}
$S_k = V_k\sum_m Y_{km}^* V_m^*$, generation minus load; with $V_k = \tilde V_k\angle\phi_k$, $Y_{km} = \abs{Y_{km}}\angle\delta_{km}$: $P_k = \sum_m \tilde V_k\tilde V_m\abs{Y_{km}}\cos(\phi_k - \phi_m - \delta_{km})$, $Q_k$ the same with $\sin$.
Bus types: PQ (load), PV (generator), slack (one; $\tilde V$, $\phi = 0$ given).
Implicit (Newton): unknowns $x = [\phi\ (\text{PQ, PV});\ \tilde V\ (\text{PQ})]$, mismatch $f(x) = [P_k(x) - P_k^{given}\ (\text{PQ, PV});\ Q_k(x) - Q_k^{given}\ (\text{PQ})]$, $x \leftarrow x - J^{-1}f(x)$ from $\tilde V = 1$, $\phi = 0$, with $J = [\partial P/\partial\phi,\ \partial P/\partial\tilde V;\ \partial Q/\partial\phi,\ \partial Q/\partial\tilde V]$.
Explicit (after): slack $P$, $Q$ and PV-bus $Q$ from their own equations.
$P_{loss} = \sum_k P_k = \sum_{\text{lines}}\abs{I_{km}}^2 R_{km}$.



\section{Power line models}
\paragraph{Models.}
% Per unit length: series $r$, $x$ [$\Omega$/km]; shunt $b$ [S/km], line charging, capacitive.
Length $\ell$: $Z = (r + jx)\ell$, $Y = jb\ell$.
Short line: $Z$ only.
$\pi$ model: $Z$ in series, $Y/2$ to neutral at each end.
% All values per phase; the line is one branch of the per-phase circuit.
% Ampacity: maximum safe $\abs{I}$.

\paragraph{Setup.}
Sending end $S$, receiving end $R$.
Load as parallel $G$ and $jB$: $Y_L = S_\phi^*/\tilde V_R^2 = (P_\phi - jQ_\phi)/\tilde V_R^2$, i.e.\ $R_L = \tilde V_R^2/P_\phi$, $X_L = \tilde V_R^2/Q_\phi$; $Y_R = Y_L + Y_{sh}/2$ ($\pi$ model), $Y_R = Y_L$ (short line).
$I_{line} = V_R Y_R$, $V_S = V_R(1 + ZY_R)$, $I_S = I_{line} + V_S Y_{sh}/2$.
$P_{loss} = 3\operatorname{Re}(V_S I_S^*) - P_{load} = 3\abs{I_{line}}^2 R$.
$VR = (\abs{V_{R,NL}} - \abs{V_{R,FL}})/\abs{V_{R,FL}}$ at fixed $V_S$; no load is $Y_L = 0$: short line $V_{R,NL} = V_S$, $\pi$ model $V_{R,NL} = V_S/(1 + ZY/2)$.




\section{Magnetism and magnetic circuits}

\paragraph{Magnetism.}
Magnetic field intensity $H$ [A/m].
Permeability $\mu = \mu_r\mu_0$ [H/m]; $\mu_0 = 4\pi\times10^{-7}$ H/m (air).
Magnetic flux density $B = \mu H$ [T $=$ Wb/m$^2$].
Magnetic flux $\Phi = BA$ [Wb] through a cross-section of area $A$ [m$^2$].
Flux linkage $\lambda = N\Phi = Li$ [Wb$\cdot$turn] for a coil of $N$ turns.
Faraday: $v = d\lambda/dt = N\,d\Phi/dt$ [V].
Ampère: $\oint H\,dl = I_{enc}$, the current passing through the loop formed by the path.
For the mean core path of length $l$ [m] with uniform $H$, and a coil of $N$ turns carrying $i$ [A] threading the loop: $Hl = Ni$.
$Ni$ is the magnetomotive force (MMF) $F$ [A$\cdot$turn].

\paragraph{Magnetic circuit.}
Electric $\iff$ Magnetic:
EMF $V$ [V] $\iff$ MMF $F = Ni$ [A$\cdot$turn];
current $I$ [A] $\iff$ flux $\Phi$ [Wb];
resistance $R$ [$\Omega$] $\iff$ reluctance $\mathcal R = l/(\mu A)$ [A/Wb], for a segment of length $l$, area $A$, permeability $\mu$;
KVL $\iff \sum F - \sum\mathcal R\Phi = 0$ around a loop (Ampère);
KCL $\iff \sum\Phi = 0$ at a junction (flux continuity);
Ohm's law $V = IR \iff F = \Phi\mathcal R$; series and parallel combine like resistance.
Coil $=$ MMF source of value $Ni$, direction by the right-hand rule.
One $\mathcal R$ per segment of distinct $l$, $\mu$, or $A$.
Air gap: $\mathcal R_{gap} = l_{gap}/(\mu_0 A)$ dominates, since $\mu_0 \ll \mu_{core}$.

\paragraph{Inductance.}
Coil of $N$ turns on a single core of length $l$, area $A$:
$L = \lambda / i = (N\Phi) / i = N^2/\mathcal R$ [H].
Several segments: $L = N^2/\mathcal R_{total}$, with $\mathcal R_{total}$ seen by the coil.

\paragraph{Saturation.}
$B$ cannot exceed $B_{sat}$ anywhere in the core.
Procedure: $F = Ni_{peak}$; $\Phi$ from the magnetic circuit; $B_k = \Phi_k/A_k$ in each segment $k$; $B_{sat} \ge \max_k B_k$, in the smallest $A_k$ carrying the largest $\Phi_k$.

\paragraph{Losses.}
Core, eddy: flux induces voltage in the core, voltage drives current, current heats; $\propto f^2B^2$; fix with laminations and silicon steel.
Core, hysteresis: area of the $B$--$H$ loop per cycle; $\propto f\times$ area; fix by material.
Winding, skin: AC current crowds to the surface; fix with Litz wire.
Winding, proximity: one conductor's field induces eddy currents in its neighbors, worse with more layers; fix with fewer or interleaved layers, Litz wire.



\section{Transformers}
Two windings on one core, $N_1$ turns on the primary and $N_2$ on the secondary.
Faraday on the primary, $V_1 = j\omega N_1 \Phi$, fixes the core flux $\Phi$.
The same $\Phi$ links the secondary, so $V_2 = j\omega N_2 \Phi$ and $V_1/N_1 = V_2/N_2$, with or without a load.
A load draws $I_2$; its MMF $N_2 I_2$ would change $\Phi$, so the primary draws the $I_1$ that cancels it.
Dot: current into the dotted terminal of either winding drives flux the same way around the core.
Define $I_1$ into the primary dot and $I_2$ out of the secondary dot; then $V_1$, $V_2$ are both $+$ at the dots and MMF balance is $N_1 I_1 - N_2 I_2 = \mathcal R\Phi$.
Autotransformer: one winding with a tap; primary and secondary share turns.
Tap changer: moves the tap to change $N_1/N_2$ for voltage control.

\paragraph{Ideal.}
$\mu\to\infty \implies \mathcal R\to 0 \implies N_1 I_1 = N_2 I_2$.
Turns ratio $a = N_1/N_2$, so $V_1 = aV_2$, $I_1 = I_2/a$, and $S_1 = V_1 I_1^* = S_2$.

\paragraph{Equivalent circuit (nonideal).}
Refer across the ideal $a{:}1$: $Z_{eq1} = a^2 Z_{eq2}$.
Circuit: primary terminals $\to Z_1 = R_1 + jX_1 \to (R_c \parallel jX_m) \to$ ideal $a{:}1 \to Z_2 = R_2 + jX_2 \to$ secondary terminals; $X_m = \omega L_m$ with $L_m = N_1^2/\mathcal R_{tot}$ (energy stored in core), $R_c$ core loss, $R_k$ winding resistance, $X_k$ leakage reactance.
$Z_{eq1} = Z_1 + a^2 Z_2$; $Z_{sh} = R_c \parallel jX_m$.
$I_1 = V_1/\big(Z_1 + Z_{sh}\parallel(a^2 Z_2 + Z_L')\big)$; with $Z_1 = Z_2 = 0$, $I_1 = V_1/Z_L' + V_1/Z_{sh}$.
% OC test (secondary open; given $\tilde V_{oc}$, $\tilde I_{oc}$, $P_{oc}$): $Q_{oc} = \sqrt{(\tilde V_{oc}\tilde I_{oc})^2 - P_{oc}^2}$, $R_c = \tilde V_{oc}^2/P_{oc}$, $X_m = \tilde V_{oc}^2/Q_{oc}$.
% SC test (secondary shorted; given $\tilde V_{sc}$, $\tilde I_{sc}$, $P_{sc}$): $R_{eq1} = P_{sc}/\tilde I_{sc}^2$, $\abs{Z_{eq1}} = \tilde V_{sc}/\tilde I_{sc}$, $X_{eq1} = \operatorname{sqrt}(\abs{Z_{eq1}}^2 - R_{eq1}^2)$.

\paragraph{Solving.}
Efficiency $\eta = P_{\rm out}/(P_{\rm out} + P_{\rm core} + P_{\rm copper})$.
$P_{\rm core} = P_{oc}$, constant at rated voltage.
$P_{\rm copper} = \tilde I_1^2 R_1 + \tilde I_2^2 R_2 = \tilde I_2^2 R_{eq2}$.
Voltage regulation $VR$: $V_{2,NL} = V_1/a$ and $V_{2,FL} = V_1/a - I_2 Z_{eq2}$ at the given load; $VR = 0$ for an ideal source.
Leading pf: the component of $I_2 Z_{eq2}$ along $V_{2,FL}$ can be negative, so $VR$ can be negative.

\paragraph{Three-phase.}
Three single-phase transformers, each side Y or $\Delta$; bank rating $S_{3\phi} = 3 S_{\rm transformer}$.
Winding quantities: Y side $V_w = V_{ln}$, $I_w = I_{\rm line}$; $\Delta$ side $V_w = V_{ll}$, $I_w = I_{\rm line}/\sqrt3$.
Per-transformer turns ratio $a = V_{1,w}/V_{2,w}$; line-to-line ratio $V_{1,ll}/V_{2,ll} = a$ (Y--Y, $\Delta$--$\Delta$), $a/\sqrt3$ ($\Delta$--Y), $\sqrt3 a$ (Y--$\Delta$).
$\Delta$--Y and Y--$\Delta$ shift voltages and currents by $30^\circ$ between sides; magnitudes and powers unaffected.
Procedure: convert line to winding quantities and $S_\phi = S_{3\phi}/3$; solve the single-phase problem on one transformer; convert back; add the $30^\circ$ only if a phasor is needed relative to the other side.




\section{Machines}
Rotor: rotating part. Stator: fixed part.
Field winding: makes the main flux from DC current $i_f$ or a permanent magnet.
Armature winding: has the EMF induced in it; its current $i_a$ carries the converted power.
% DC: field on the stator, armature on the rotor. 
% Synchronous: field on the rotor, armature on the stator.
% Motor: electrical in, mechanical out. Generator: the reverse.
$E_{\rm electric} = E_{\rm mech} + E_{\rm field} + E_{\rm loss}$.
Four facts: current in a wire makes a field (transformer primary); a wire carrying $i$ in field $B$ feels $\vec F = i\,\vec l\times\vec B$ (motor); a changing field makes a voltage in a coil (transformer secondary); a wire moving at $\vec v$ in field $B$ has EMF $e = (\vec v\times\vec B)\cdot\vec l$ (generator).
Left-hand rule, motor: pointer $B$, middle $i$, thumb $F$.

\paragraph{DC machines.}
Stator: field, fixed flux from a permanent magnet or DC field winding. 
Rotor: armature, fed DC through brushes; the commutator reverses each coil's current to keep one diretion of torque.
(+) cheap drive, stable. ($-$) commutator wear.

Circuit: terminals $\to R_a \to L_a \to$ back EMF $E = k_e\omega$, so $V = R_a i_a + L_a\,di_a/dt + k_e\omega$.
Steady state drops $L_a$ term.
Motor ($i_a \to R_a\to E$): $V = R_a i_a + k_e\omega$. 
Generator ($E\to R_a\to i_a$): $V = k_e\omega - R_a i_a$.
Back EMF constant $k_e = (V - i_a R_a)/\omega$ [V$\cdot$s/rad], from a no-load test. Torque constant $k_t$ [N$\cdot$m/A] in $T = k_t i_a$.
Power balance $P_{\rm mech} = E i_a = T\omega \implies k_e = k_t$ and $i_a = T/k_t$.
$P_{\rm in} = Vi_a$.
Torque--speed: $\omega = V/k_e - TR_a/(k_e k_t)$, a line from no-load speed $\omega_0 = V/k_e$ ($T = 0$) to stall torque $T_0 = k_t V/R_a$ ($\omega = 0$).


\paragraph{Synchronous machines.}
Stator: armature, three-phase winding; its currents make a field rotating at $\omega_s$ and carry the AC power.
Rotor: field winding with DC $i_f$ through slip rings (or a permanent magnet); its flux locks to the rotating field, so the shaft turns at $\omega_s$.
$p$ pole pairs: $\omega_s = \omega_e/p = 2\pi f_e/p$ [rad/s], since the stator field advances one pole pair, $1/p$ of a turn, per electrical cycle.
At standstill the torque alternates at $f_e$ and averages zero (solution: squirrel cage or ramp $f_e$)
(+) pf set by $i_f$; speed fixed by $f_e$. ($-$) not self-starting; needs a DC field supply.
% Magnetic circuit: $(N_f i_f \to \mathcal R_{\rm gap} \to \mathcal R_{\rm rotor} \to \mathcal R_{\rm gap} \to N_f i_f) \parallel (\mathcal R_{\rm stator} \parallel \mathcal R_{\rm stator})$.
Per-phase circuit: EMF $E_a \to R_s \to jX_s \to$ terminal $V_a$; generator has $I_a$ from machine to grid so $(E_a - V_a) = (R_s + jX_s)I_a$; reverse direction for motor.
Typically set $R_s = 0$.
$E_a = k_f i_f\omega_e\angle\delta = (L_f/\sqrt2)\, i_f\omega_e\angle\delta$; $\delta = \phi_{E_a} - \phi_{V_a}$ for a generator; reverse for motor so $\delta$ always $> 0$.
$S_{3\phi} = 3V_aI_a^* = 3V_a((E_a - V_a)/jX_s)^*$ with $\tilde V_a\angle 0$ and $\tilde E_a\angle \delta$.
% \paragraph{Solving SMs.}
% Given $|S_{3\phi}|$, pf, $\tilde V_{ll}$, $X_s$: $I_a = (|S_{3\phi}|/(3\tilde V_a))\angle(-\theta)$ with $\theta$ from pf; then $E_a = V_a + jX_s I_a$ (generator) gives $\tilde E_a$, so $i_f$, and $\delta$.
% Procedure: given $\tilde E_a$, $P_{3\phi}$, $\tilde V_{ll}$, $X_s$, get $\delta$ from the $P$ equation; then $I_a = (E_a - V_a)/(jX_s)$ and $Q_{3\phi}$.
Torque--speed: vertical line at $\omega_s$.
Torque--$\delta$: $T = P_{3\phi}/\omega_s = 3\tilde V_a\tilde E_a\sin\delta/(X_s\omega_s)$, a sine, max at $\delta = 90^\circ$, unstable past it.
$T = (\pi/2)\,p\,\Phi_r F_f\sin\delta$ with air-gap flux per pole $\Phi_r$ and field MMF $F_f = N_f i_f$; raise $T$ by $\Phi_r$ (limited by $B_{\rm sat}$) or $F_f$ (limited by field heating).



\end{document}